\documentclass[10pt]{article}

\usepackage[utf8]{inputenc}
\usepackage{floatrow}

\usepackage{algorithm, algpseudocode}
\let\oldReturn\Return
\renewcommand{\Return}{\State\oldReturn}
\newcommand{\N}{\mathbb{N}}
\newcommand{\R}{\mathbb{R}}
\usepackage[T1]{fontenc}
\usepackage{enumitem}
\usepackage{hyperref}
\usepackage{scrextend}
\usepackage{amsmath}
\usepackage{amsfonts}
\usepackage{stmaryrd}
\usepackage{graphicx}
\usepackage{color}
\usepackage{listings}
\usepackage{wrapfig}
\usepackage[hmargin=1.25in,vmargin=1.25in]{geometry}

%title setup
\title{TD1 : l'ordonnanceur du noyau Linux}
\author{Romain PEREIRA}
\date{25/10/2018}

% table of contents setup
\renewcommand{\contentsname}{Sommaire}
\usepackage{etoolbox}
\patchcmd{\thebibliography}{\section*{\refname}}{}{}{}

\usepackage[utf8]{inputenc}
\usepackage[T1]{fontenc}
\usepackage[frenchb]{babel}

\setlength{\parindent}{0cm}
\setlength{\parskip}{1ex plus 0.5ex minus 0.2ex}
\newcommand{\hsp}{\hspace{20pt}}
\newcommand{\HRule}{\rule{\linewidth}{0.5mm}}

\hypersetup{
    colorlinks,
    citecolor=black,
    filecolor=black,
    linkcolor=blue,
    urlcolor=red
}

\lstset{language=C,
                basicstyle=\ttfamily,
                keywordstyle=\color{blue}\ttfamily,
                stringstyle=\color{red}\ttfamily,
                commentstyle=\color{cyan}\ttfamily,
                morecomment=[l][\color{magenta}]{\#}
}

\lstdefinelanguage{Markdown}{}

\begin{document}
    
    \begin{titlepage}
        \begin{sffamily}
            \begin{center}
                
                                
                \begin{figure}[h!]
                    \includegraphics[width=6cm]{ensiie.jpeg}
                \end{figure}

                \HRule \\[0.4cm]
                { \huge \bfseries Architecture d'un système d'exploitation\\[0.4cm] }
                \HRule \\[2.0cm]
                
                { \huge \bfseries TD2\\[0.5cm] }

                \textsc{\Large L'ordonnanceur du noyau Linux}\\[2.0cm]

                \vfill
                \begin{minipage}{0.4\textwidth}
                    \begin{flushleft} \large
                        \emph{Etudiant:} Romain \textsc{Pereira}\\
                    \end{flushleft}
                \end{minipage}
                \begin{minipage}{0.4\textwidth}
                    \begin{flushright} \large
                        \emph{Encadrants:}  M. \textsc{Loussert}\\
                                            M. \textsc{Taboada}
                    \end{flushright}
                \end{minipage}
                \\[2.0cm]
                {\large 25/10/2018}
            \end{center}
        \end{sffamily}
    \end{titlepage}
    
    \maketitle
    \tableofcontents
    \section{Préambule}
    
    Le but de ce TP est de comprendre le fonctionnement d'un ordonnanceur, en étudiant l'ordonnanceur du noyau Linux.

    \newpage
    \section{Fonctionnement de l'ordonnanceur}
        \paragraph{Q.1}
        Nos ordinateurs de bureau dispose généralement d'un seul processeur (avec 1, 2, 4 ou 8 coeurs d'execution).
        
        Lorsque plusieurs processus s'execute sur la machine, il faut donc qu'il
        y ait un méchanisme de partage du processeur : il est le seul point d'execution pour les processus.
    
        C'est la fonction principale de l'ordonnanceur.
        Il choisit dans quel ordre plusieurs processus peuvent être executé sur un même processeur.
        
        \textbf{N.B} : l'ordonnanceur est lui même un processus.
                
        L'execution de l'ordonnanceur peut avoir lieu pour répondre aux problèmes suivants:
        
        - Lorsqu'un processus se clone (\textit{fork()}) : lequel du père ou du fils doit avoir la priorité d'éxecution?
        
        - Lorsque le processus courant s'arrête : quel processus doit prendre la main ensuite
        
        - Lorsqu'un processus bloque sur une entrée/sortie, ou qu'il est en attente : doit t'il continuer de bloquer le processeur alors qu'il ne travaille pas?
        
        - Lorsqu'une entrée/sortie est interrompu : peut être que le processus a fini son job, ou bien c'est un signal indiquant qu'un processus précèdemment en attente entrée/sortie peut maintenant effectué son travail.
        
        
        Aussi, on distingue 2 grandes politiques décidant de l'execution de l'ordonnanceur.

        - non-pre-emptive : l'ordonnanceur ne tourne que si le processus utilisateur s'arrête, bloque, ou demande explicitement à changer de processus.
        Si le processus ne rends pas la main pendant 100 ans, l'ordonanneur attendra impassiblement pendant 100 ans.

        \label{preempt}
        - pre-emptive : Chaque processus à un temps d'execution fixe.
        L'ordonnanceur est executé à intervalle régulier, et détermine le prochain prochain processus qui doit être executé sur le prochaine 'créneau'.


        \paragraph{Q.2}
        La fonction \textit{sched\_yield} passe la main sur le processeur d'un processus à un autre
        (ou d'un processus vers lui même s'il reste le plus 'prioritaire').
        
        S'il n'y a pas d'autre processus devant être executer, cette fonction s'arrête.
        
        \paragraph{Q.3}
        Code modifié
        \lstset{language=C}
\begin{lstlisting}[frame=single]
[...]
printk("current=%p\n", current);
schedule();
printk("current=%p\n", current);
[...]
\end{lstlisting}

        Compilation de l'image linux modifié et déploiement
\lstset{language=bash}
\begin{lstlisting}[frame=single]
cd ~/debian_kernel/linux-4.9.30/
make bindeb-pkg
dpkg -i ../linux-image-4.9.30_4.9.30-2_amd64.deb
reboot
\end{lstlisting}

Programme utilisateur 'main.c'
\lstset{language=C}
\begin{lstlisting}[frame=single]
int main() {
    sched_yield();
    return 0;
}
\end{lstlisting}

Compile et lance le programme utilisateur
\begin{lstlisting}[frame=single]
gcc main.c
./a.out
dmesg
\end{lstlisting}

Sortie dmesg
\begin{lstlisting}[frame=single]
[  170.029953] current=ffff9337bbe47140
[  170.029956] current=ffff9337bbe47140
\end{lstlisting}

=> Le pointeur 'current' reste inchangé
Ceci est dû au fait que peu de processus tourne sur la machine virtuelle.
Le processus prioritaire reste le même entre 2 appels de 'schedule()'

En lançant d'autres processus sur la machine, la valeur du pointeur aurait changé


\paragraph{Q.4}
On génère d'abord les 'ctags' pour naviguer plus facilement dans le code source.

\lstset{language=bash}
\begin{lstlisting}[frame=single]
cd ~/debian_kernel/linux-4.9.30/
ctags -R .
\end{lstlisting}


Lecture de 'core.c'
\lstset{language=C}
\begin{lstlisting}[frame=single]
ligne  profondeur/schedule()  execution
4898  0  schedule();

3454  1  sched_submit_work(task);
3438  2  ...

//debut boucle : desactive le préemption
3456  1  preempt_disable();

3457  1  __schedule(false);
3342  2  cpu_rq(); // recuperes la 'run queue' du CPU
3343  2  prev = rq->curr; // enregistre le processus courant dans la run queue
3391  2  next = pick_next_task(rq, prev, cookie); // recuperes la task suivante

// réactive la préemption
3457  1  sched_preempt_enable_no_resched();

3459  1  need_resched();
// tant que 'need_resched()' renvoit 'vrai',
// on boucle sur partir de preempt_disable()
\end{lstlisting}
        
\paragraph{Q.5} La fonction principal de l'ordonnanceur est 'schedule()'

\lstset{language=C}
\begin{lstlisting}[frame=single]
// choisit le prochain processus à être executé
// avant l'appel de cette fonction, un processus tourne
schedule();
// apres l'appel, un autre processus peut tourner
\end{lstlisting}
        
\paragraph{Q.6}
La pré-emption est un mecanisme d'ordonnancement qui consiste à donner un temps maximal d'execution à un processus,
au dela duquelle l'algorithme d'ordonnancement sera relancé (c.f \ref{preempt} )

Si la préemption est active lors de l'execution de l'ordonnanceur, ...
 
 \paragraph{Q.7}
 La 'run-queue' est une structure de données (de type tableau). Elle liste la priorité d'execution des processus.
 A chaque processeur est associé une run-queue.
 
 La run-queue est verouillé au début de la fonction \textiŧ{schedule} (alors que le processus 'prev' vient de finir son créneau d'execution)
 
 Ensuite, le prochain processus devant être executé est élu par l'ordonnanceur : c'est le processus 'next'
 
 Une fois l'election terminé, le verrou de la run-queue est retiré et le processus 'next' commence son execution.
 
 Ce verrou sur la run-queue permet donc de ne pas modifier la priorité des processus pendant l'éléction.
 Une fois l'election terminé, le processus 'next' libère ce verrou pour que la priorité des processus puissent être mis à jour par la suite.
 
\paragraph{Q.8} ligne 3258 core.c
\lstset{language=C}
\begin{lstlisting}[frame=single]
struct task_struct *pick_next_task (struct rq *rq,
				    struct task_struct *prev,
				    struct pin_cookie cookie);
\end{lstlisting}
 
 \paragraph{Q.9}
 La fonction \textit{context_switch} se trouve à la ligne 2862 du fichier 'core.c'.
 Elle est declaré en '__always_inline', ce qui signifie que le code de cette fonction doit être duppliqué par le compilateur,
 et ré-insérer dans le binaire executable à chaque appel vers cette fonction (voir \ref{inlinefunc} )
 
 Cette fonction rétablie l'état du processus sur le processeur, état dans lequel il avait été stoppé lors de son dernier créneau d'execution (ou dans un nouvel état si c'est un nouveau processus).
 
 
 
 
 \paragraph{Q.10} TODO

    \newpage
    \section{Références}
    \begin{thebibliography}{}

        \bibitem{modernos}\label{modernos}
        Modern Operating Systems - Andrew S. Tanenbaum, Herbert Bos (page 149-166)\newline
        \href{https://github.com/concerttttt/books/blob/master/Modern Operating Systems 4th Edition--Andrew Tanenbaum.pdf}
        {\textit{https://github.com/concerttttt/books/}}\newline
	
        \bibitem{inlinefunc}\label{inlinefunc}
        Inlining in Linux - Linux Kernel Documentation\newline
        \href{https://www.kernel.org/doc/local/inline.html}
        {\textit{https://www.kernel.org/doc/local/inline.html}}\newline
        
    \end{thebibliography}
    
\end{document}
